\section*{A043: genuine disjoint prime gates}

All Boolean functions in this note take values in $\{-1,1\}$.  Before
factorization, delete every inessential variable.  The resulting function is
called its \emph{essential core}.

\paragraph{Honest disjoint factorization.}
Let $f$ have essential variable set $X$ of size $n$.  An honest nontrivial
disjoint factorization is an identity
\[
 f(x)=G\bigl(g_1(x_{B_1}),\ldots,g_k(x_{B_k})\bigr),             \tag{1}
\]
where $(B_1,\ldots,B_k)$ is a partition of $X$, $2\le k<n$, every
$g_i$ is nonconstant and depends essentially on every variable in $B_i$,
and $G$ depends essentially on all $k$ input ports.  The restriction $k<n$
excludes the tautological representation using $n$ literals.

An essential function of arity at least two is \emph{disjoint-prime} if it
has no factorization (1).  Constants are nullary atoms and signed literals
are unary atoms; neither is called prime.  Prime classes are taken modulo
permutation of variables and input/output negations: an input negation is
absorbed into the adjacent child and an output negation into the parent.

Affine parity follows the same convention.  Parity on $n\ge3$ variables is
composite, since it factors through a binary XOR/XNOR gate and parity on the
remaining variables.  Binary XOR and XNOR are prime only because arity two
is the base arity.  One may flatten consecutive parity nodes for display,
but this is not a different factorization notion.

\begin{theorem}[Existence and termination]
Every finite Boolean function has a rooted composition tree whose leaves are
constants or signed literals and whose internal nodes are disjoint-prime
gates.  At each internal node, its child subtrees use pairwise disjoint sets
of the original variables.  Every root-to-leaf path has length at most
$n-1$.
\end{theorem}

\begin{proof}
Pass to the essential core.  Constants and literals are atoms.  If an
essential function of arity $m\ge2$ is prime, stop.  Otherwise choose an
honest factorization (1), and recursively factor both the outer gate $G$ and
every non-atomic child $g_i$.

Every recursive successor has strictly smaller arity than $m$: the outer
gate has arity $k<m$, and every child block has size at most $m-k+1<m$.
Thus arity strictly decreases along every branch.  Since branching is finite,
the recursion terminates, and its depth is at most $n-1$.  Disjointness is
preserved because factoring a child only partitions that child's block, while
factoring an outer gate partitions its already-disjoint input ports.
\end{proof}

No uniqueness is asserted.  Associative gates, especially parity, AND and
OR, have many valid trees.  The theorem supplies existence and a terminating
induction parameter, which is all that the entropy implication needs.

\begin{proposition}[A structural restriction on genuine primes]
Every essential disjoint-prime Boolean gate of arity $n\ge3$ is
non-canalizing: neither section in any coordinate is constant.
\end{proposition}

\begin{proof}
Suppose the section $f|_{x_i=a}$ is the constant $c$.  Let
$g(x_{X\setminus\{i\}})=f|_{x_i=1-a}$.  Essentiality of $f$ implies that
$g$ depends on every remaining variable.  Define the binary gate $G$ by
$G(a,v)=c$ and $G(1-a,v)=v$.  Then
\[
 f(x)=G(x_i,g(x_{X\setminus\{i\}})).
\]
Both inputs of $G$ are essential, and $2<n$, so this is an honest
factorization, contradicting primality.
\end{proof}

Thus all essential canalizing gates of arity at least three, including the
usual iterated AND/OR and selector forms, are removed before the prime
obligation.  This restriction does not solve the remaining class: arbitrary
non-canalizing gates can still be prime.
